%!TEX root = ./main.tex

\section{Wrap-up}

% SECTION TIMING: 69:00--75:00. Use the comparison only to close the story, then collect the exit ticket.

% TEACHING NOTE (2:30): Read down by generation, then across one row. Do not test acronym memorization.
\begin{frame}{Five generations in one low-resolution map}
  \vspace{0.3em}
  \centering
  \scriptsize
  \setlength{\tabcolsep}{3pt}
  \renewcommand{\arraystretch}{1.32}
  \begin{tabular}{>{\bfseries\raggedright\arraybackslash}p{0.45cm}
      >{\raggedright\arraybackslash}p{1.45cm}
      >{\raggedright\arraybackslash}p{1.55cm}
      >{\raggedright\arraybackslash}p{2.30cm}
      >{\raggedright\arraybackslash}p{1.30cm}
      >{\raggedright\arraybackslash}p{2.40cm}}
    \toprule
    & \textbf{Radio} & \textbf{RAN control} & \textbf{Core} & \textbf{Mode} & \textbf{Example} \\
    \midrule
    1G & Base station & MTSO & Analog core & Circuit & Voice in a car \\
    2G & BTS & BSC & MSC + HLR/VLR & Circuit; GPRS packet & SMS, roaming \\
    3G & NodeB & RNC & CS + PS cores & Both & Mobile web / video \\
    4G & eNodeB & In the eNodeB & EPC, all-IP & Packet & Maps, apps, streaming \\
    5G & gNodeB & CU (RRC), split CU/DU/RU & Service-based 5GC & Packet + QoS flows & Flexible services + edge \\
    \bottomrule
  \end{tabular}

  \glossline{CS = circuit-switched \quad PS = packet-switched \quad
    QoS = Quality of Service\\
    EPC = Evolved Packet Core \quad 5GC = 5G Core}

  \vspace{0.5em}
  \takeaway{The radio name changes.\\The deeper story is where control lives.}
\end{frame}

% TEACHING NOTE (2:00): Point along the path and ask students to finish each phrase aloud.
\begin{frame}{Follow the same action one last time}
  \vspace{0.4em}
  \centering
  \begin{tikzpicture}[node distance=0.2cm]
    \node[netnode,minimum width=1.58cm] (g1) {1G\\[-1pt]{\scriptsize circuit}};
    \node[netnode,minimum width=1.58cm,right=of g1] (g2) {2G\\[-1pt]{\scriptsize identity}};
    \node[netnode,minimum width=1.58cm,right=of g2] (g3) {3G\\[-1pt]{\scriptsize packet data}};
    \node[netnode,minimum width=1.58cm,right=of g3] (g4) {4G\\[-1pt]{\scriptsize all-IP}};
    \node[corenode,minimum width=1.58cm,right=of g4] (g5) {5G\\[-1pt]{\scriptsize modular core}};
    \node[oranode,minimum width=1.58cm,right=of g5] (o) {O-RAN\\[-1pt]{\scriptsize open control}};
    \draw[flow] (g1)--(g2); \draw[flow] (g2)--(g3); \draw[flow] (g3)--(g4);
    \draw[flow] (g4)--(g5); \draw[flow] (g5)--(o);
  \end{tikzpicture}

  \vspace{1.2em}
  \Large
  The packet still needs \textcolor{ranblue}{radio access},\\[0.25em]
  \textcolor{coreorange}{identity + policy + routing},\\[0.25em]
  and now sometimes \textcolor{ranpurple}{programmable control}.

  \vspace{1.0em}
  \takeaway{Each generation relocates functions.\\O-RAN changes who may control them.}
\end{frame}

% TEACHING NOTE (2:30): Collect on paper or LMS. Require one boundary and one sentence.
\begin{frame}{Exit ticket: draw the boundaries}
  \vspace{0.6em}
  \begin{center}
  \begin{tikzpicture}[node distance=0.32cm]
    \node[netnode,minimum width=1.55cm] (ue) {UE};
    \node[netnode,right=of ue,minimum width=2.25cm] (ran) {Radio system};
    \node[corenode,right=of ran,minimum width=2.35cm] (core) {Network services};
    \node[netnode,right=of core,minimum width=2.35cm,font=\small\bfseries] (data) {PSTN / Internet};
    \draw[flow] (ue)--(ran); \draw[flow] (ran)--(core); \draw[flow] (core)--(data);
  \end{tikzpicture}
  \end{center}

  \vspace{1.0em}
  \begin{enumerate}\itemsep6pt
    \item Mark the \textcolor{ranblue}{RAN} and \textcolor{coreorange}{Core} boundary.
    \item Add one \textcolor{ranpurple}{O-RAN} interface or controller.
    \item Answer in one sentence:\\
      \textbf{What changes next---radio, core, or control?}
  \end{enumerate}
\end{frame}
